\documentclass[12pt]{article}
\usepackage[T1]{fontenc}
\usepackage[utf8]{inputenc}
\usepackage{lmodern}
\usepackage{textcomp}
\usepackage{enumitem}
\usepackage{geometry}
\geometry{margin=1in}
\begin{document}
\begin{center}
Answer Key - History - Graduate Level Assessment 20 \
\small (Answers are ordered exactly as in the question paper.)
\end{center}

\section*{Multiple Choice}
\begin{enumerate}[label=\arabic*.]
\item \textbf{Answer:} B
\\ \textit{Options:} A) A forest fire burned away the dense vegetation that concealed the cemetery.; B) Thieves raiding the tombs were caught by the police.; C) Archaeologists discovered an exact replica of the city in the map room.; D) The cemeteries were revealed after a massive earthquake.; E) A flock of sheep that had been living there fell into a tomb.
\\ \textit{Note:} The discovery of the royal cemeteries of the Moche was primarily due to the illegal activities of tomb raiders, whose capture by police led to the awareness and subsequent investigation of these significant archaeological sites.
\item \textbf{Answer:} A
\\ \textit{Options:} A) use of a calendar system based on celestial bodies; B) practice of mummification for the dead; C) ritual consumption of maize beer; D) construction of underground dwellings; E) ritual human sacrifice
\\ \textit{Note:} The correct answer is supported by archaeological findings that indicate the Hohokam civilization adopted a calendar system influenced by the Mesoamerican cultures, utilizing celestial bodies to mark agricultural cycles and significant events.
\item \textbf{Answer:} C
\\ \textit{Options:} A) 100; B) 500; C) 1,000; D) 5,000
\\ \textit{Note:} The correct answer is 1,000 because archaeological evidence indicates that the palace of Knossos, a significant center of Minoan civilization on Crete, featured an intricate layout with approximately 1,000 rooms, reflecting its complexity and the advanced architectural skills of the time.
\item \textbf{Answer:} A
\\ \textit{Options:} A) 5100 B.P.; B) 6200 B.P.; C) 7000 B.P.; D) 4000 B.P.; E) 1200 B.P.
\\ \textit{Note:} The correct answer, 5100 B.P., corresponds to the earliest archaeological findings that indicate the use of metals, particularly in the context of ancient Andean civilizations, highlighting a significant development in metallurgy in South America during that period.
\item \textbf{Answer:} E
\\ \textit{Options:} A) a vast network of roads and trade routes.; B) the expansion and control of waterways for transportation and irrigation.; C) the cultivation and trade of cacao and maize.; D) religious beliefs that required extensive and escalating human sacrifice.; E) tribute in the form of gold, jade, feathers, cloth, and jewels.
\\ \textit{Note:} The Aztec Empire's expansion relied on a tribute system where conquered peoples were required to supply valuable resources such as gold, jade, feathers, cloth, and jewels, which not only enriched the empire but also reinforced its political dominance and control over vast territories.
\end{enumerate}

\section*{True / False}
\begin{enumerate}[label=\arabic*.]
\setcounter{enumi}{5}
\item \textbf{Answer:} True
\\ \textit{Options:} A) True; B) False
\\ \textit{Note:} According to the passage: 1918 to 1920
\item \textbf{Answer:} True
\\ \textit{Options:} A) True; B) False
\\ \textit{Note:} According to the passage: Treaty of Vienna
\item \textbf{Answer:} False
\\ \textit{Options:} A) True; B) False
\\ \textit{Note:} The correct answer is: 1676
\item \textbf{Answer:} False
\\ \textit{Options:} A) True; B) False
\\ \textit{Note:} The correct answer is: September 1829
\item \textbf{Answer:} False
\\ \textit{Options:} A) True; B) False
\\ \textit{Note:} The correct answer is: Garmisch-Partenkirchen, Germany
\end{enumerate}

\section*{Long Form Response}
\begin{enumerate}[label=\arabic*.]
\setcounter{enumi}{10}
\item \textbf{Answer:} half-hourly
\\ \textit{Note:} Expected answer: half-hourly
\item \textbf{Answer:} 1987
\\ \textit{Note:} Expected answer: 1987
\end{enumerate}

\vfill
\noindent\textit{End of Answer Key}
\end{document}